\chapter{Shock waves}
\label{ch:shocks}

\tocome{Drawn from the course's Units~3 to~6. The shock of finite thickness,
steady in its own frame: the Reynolds transport theorem of
\cref{sec:balance} across it gives the jump conditions whatever its
structure, and the Rayleigh line and the Hugoniot; the energy split in
two; the entropy made inside the front, fixed by the end states while its
mechanism sets only the thickness; the shock speed and the shock impedance
in place of $c$ and $Z$; impedance matching at a contact. And the step to
the numerics: a method may choose its own shock thickness, a few cells
wide, as long as it keeps the jump conditions, which is what method AV's
viscosity does (\cref{ch:method-av}).}
